\chapter{Robustness and tissue-generalisation analyses}\label{ch:robustness}

The results of Chapter~\ref{ch:results} come from the frozen primary benchmark. A separate
and clearly labelled layer of analyses tests how durable its small advantage is, without
altering it (Table~\ref{tab:robustness}). Three of the four planned analyses were executed;
the negative controls were configured but not run.

\begin{table}[tbp]
  \caption{The exploratory and robustness layer: questions, runs and execution status.}
  \label{tab:robustness}
  \footnotesize
  \begin{tabular}{@{}L{0.20\linewidth}L{0.41\linewidth}L{0.13\linewidth}L{0.19\linewidth}@{}}
    \toprule
    \thead{Analysis} & \thead{Question} & \thead{Run} & \thead{Status} \\
    \midrule
    Random forest M3 & Does a different model family extract more signal from the same
      features and splits? & \path{exploratory_m3_v1} & Completed, 200/200 tasks \\
    Leave-one-tissue-out & Does an expression-only model transfer to a tissue it has never
      seen? & \path{loto_v1} & Completed, 18/18 tasks \\
    Outcome-scale sensitivity & Does the advantage survive the neglog and logit scales, or
      is it a raw-scale artefact? & \path{scales_v1} & Completed, 400/400 tasks \\
    Negative controls (1,000 draws of each null) & Is the advantage larger than the pipeline
      produces when expression cannot carry signal? Are any gene sets more important than
      matched random sets? & \path{controls_v1}
      & Configured, about 10--20~h at six workers; \textbf{not executed} \\
    \bottomrule
  \end{tabular}
\end{table}

\section{A different model family does not do better}\label{sec:m3}

The random forest M3, fitted on the same splits with the same nested tuning, reached MAE
0.13148 and out-of-sample $R^2$ $+0.00643$, against 0.13040 and $+0.00935$ for M2
(Table~\ref{tab:r1}). Its gain over the tissue-only baseline, $\Delta\mathrm{MAE}=0.00275$,
is \textbf{smaller} than the 0.00383 of M2. \textbf{M3 therefore performs worse than M2, and
that is the result reported.} Within the same run, M0, M1 and M2 reproduced the primary run
to the last decimal place, which also confirms that the frozen benchmark is reproducible.

\begin{table}[tbp]
  \caption[Exploratory model comparison including the random forest M3]{Exploratory model
    comparison including the random forest M3, on the raw mSF2 scale (MAE and RMSE in mSF2
    units). M0, M1 and M2 reproduce Table~\ref{tab:performance}. Source:
    \path{results/tables/exploratory_m3_v1/model_metrics.csv}; the paired differences
    quoted in the text are from \path{comparisons.csv} in the same folder.}
  \label{tab:r1}
  \small
  \begin{tabular}{@{}l S[table-format=1.6] S[table-format=1.6] S[table-format=-1.6]
      S[table-format=-1.2] S[table-format=1.2]@{}}
    \toprule
    \thead{Model} & {\thead{MAE}} & {\thead{RMSE}} & {\thead{$R^2_{\mathrm{oos}}$}}
      & {\thead{Calibration slope}} & {\thead{Calibration intercept}} \\
    \midrule
    M0 & 0.133624 & 0.170755 & -0.009497 & -4.00 & 2.92 \\
    M1 & 0.134228 & 0.171494 & -0.018275 & -1.21 & 1.29 \\
    M2 & 0.130397 & 0.169115 & 0.009347 & 0.65 & 0.21 \\
    M3 & 0.131480 & 0.169394 & 0.006434 & 0.55 & 0.26 \\
    \bottomrule
  \end{tabular}
\end{table}

\section{Transfer to an unseen tissue is weak and conditional}\label{sec:loto}

Holding out each of the nine well-supported tissues in turn (185 components, 187 records),
the expression-only model M2X reached MAE 0.12267 and out-of-sample $R^2$ \textbf{+0.0308},
against 0.12753 and $-0.0084$ for the global mean ($\Delta\mathrm{MAE}=0.00487$ mSF2 units).
This is the only analysis in the project in which out-of-sample $R^2$ is clearly above zero.
It is, however, a different and easier comparison than the primary one --- expression alone
against a global mean, with tissue removed from both sides --- so it is not comparable with
the primary 0.00383 and does not revise it.

The pooled figure conceals an uneven picture (Table~\ref{tab:r2}). The gain is largest when
Ovary/Fallopian Tube ($+0.0156$) or Head and Neck ($+0.0153$) is held out, followed by
Breast ($+0.0119$), Bowel ($+0.0092$), Bladder/Urinary Tract ($+0.0079$), Lung ($+0.0049$)
and CNS/Brain ($+0.0025$). It is \textbf{negative for two tissues, Pancreas ($-0.0142$) and
Esophagus/Stomach ($-0.0030$)}, where the expression-only model predicted worse than the
global mean. The Pancreas loss is nearly three times the pooled gain, and Pancreas is also
the one well-supported tissue in which M2 was worse than M1 in the primary benchmark
(Table~\ref{tab:tissue}). Transfer of expression information to an unseen tissue is
therefore conditional on the tissue, and the tissue-level results, not only the pooled
gain, belong in any statement about generalisation.

\begin{table}[tbp]
  \caption[Leave-one-tissue-out performance by held-out tissue]{Leave-one-tissue-out
    performance of the expression-only model M2X by held-out tissue. $\Delta$MAE is
    MAE(global mean)$\,-\,$MAE(M2X) in mSF2 units, so positive values favour expression;
    the Spearman correlation is that of M2X within the held-out tissue. The pooled row is
    from \path{comparisons.csv} and \path{model_metrics.csv} of the same run. Source:
    \path{results/tables/loto_v1/held_out_tissue.csv}.}
  \label{tab:r2}
  \small
  \begin{tabular}{@{}l S[table-format=3.0] S[table-format=-1.4] S[table-format=-1.4]
      S[table-format=-1.3]@{}}
    \toprule
    \thead{Held-out tissue} & {\thead{Components}} & {\thead{$\Delta$MAE}}
      & {\thead{Mean prediction shift}} & {\thead{Spearman $\rho$}} \\
    \midrule
    Bladder/Urinary Tract & 11 & 0.0079 & -0.0425 & 0.327 \\
    Bowel & 21 & 0.0092 & 0.0096 & 0.231 \\
    Breast & 17 & 0.0119 & 0.0290 & 0.324 \\
    CNS/Brain & 23 & 0.0025 & 0.0735 & 0.359 \\
    Esophagus/Stomach & 23 & \bfseries -0.0030 & 0.0273 & -0.124 \\
    Head and Neck & 18 & 0.0153 & -0.0284 & 0.187 \\
    Lung & 46 & 0.0049 & -0.0221 & 0.178 \\
    Ovary/Fallopian Tube & 10 & 0.0156 & -0.0452 & 0.576 \\
    Pancreas & 16 & \bfseries -0.0142 & 0.0115 & -0.062 \\
    \midrule
    Pooled (nine tissues) & 185 & 0.0049 & {\na} & {\na} \\
    \bottomrule
  \end{tabular}
\end{table}

\section{The advantage depends on the outcome scale}\label{sec:scales}

The most important negative finding of the robustness layer concerns the outcome scale.
Repeating the M2-versus-M1 comparison with 20 repeats on two transformed scales gave
$\Delta\mathrm{MAE}=0.00023$ on the neglog scale --- essentially nothing, with the
out-of-sample $R^2$ of M2 falling to $-0.0322$ --- and 0.00980 on the logit scale
(Table~\ref{tab:r3}). The direction favours M2 on all three scales. The magnitudes, however,
are expressed in the units of each transformed scale, not in mSF2, so they cannot be
compared directly with the raw-scale 0.00383; each must be read against the baseline error
on its own scale, which is also given in the table. On the neglog scale the advantage is
negligible; on the logit scale the out-of-sample $R^2$ of M2 is slightly positive (0.0045)
while that of M1 is negative. The raw scale
remains the protocol's prespecified primary scale, and these results are reported beside it
rather than replacing it. Together they mean that the primary estimate should be read as
\textbf{small, consistent in direction and sensitive to how the outcome is expressed}.

\begin{table}[tbp]
  \caption[Outcome-scale sensitivity of the primary comparison]{Outcome-scale sensitivity of
    the primary comparison (M2 against M1). $\Delta$MAE and MAE are in the units of each
    scale and are not comparable across scales. The raw row is the primary run (10 repeats);
    the transformed scales used 20 repeats on the same 237 components and 241 records.
    Descriptive only; no confirmatory method was adopted. Sources:
    \path{results/tables/scales_v1/comparisons.csv} and \path{model_metrics.csv};
    raw scale from Table~\ref{tab:performance}.}
  \label{tab:r3}
  \small
  \begin{tabular}{@{}l S[table-format=1.5] S[table-format=1.5] S[table-format=2.0]
      S[table-format=1.4] S[table-format=-1.4] S[table-format=-1.4]@{}}
    \toprule
    \thead{Scale} & {\thead{$\Delta$MAE}} & {\thead{Record-weighted}} & {\thead{Repeats}}
      & {\thead{MAE M1}} & {\thead{$R^2_{\mathrm{oos}}$ M1}} & {\thead{$R^2_{\mathrm{oos}}$ M2}} \\
    \midrule
    raw (primary) & 0.00383 & 0.00382 & 10 & 0.1342 & -0.0183 & 0.0093 \\
    neglog & 0.00023 & 0.00017 & 20 & 0.2596 & -0.0164 & -0.0322 \\
    logit & 0.00980 & 0.00957 & 20 & 0.6398 & -0.0142 & 0.0045 \\
    \bottomrule
  \end{tabular}
\end{table}

\section{Rules governing the exploratory layer}\label{sec:rules}

Four rules keep this layer separate from the primary benchmark.

\paragraph{It cannot promote a new primary model.} M3 was configured \emph{after} the
primary results were seen, so M2 remains the primary benchmark whatever M3 scores. In the
event M3 scored worse, which is the expected outcome for a tree ensemble on 237 components
and 50 coarse features and suggests that the ceiling is set by the data rather than by the
model family.

\paragraph{Consistency with the frozen run is verified mechanically.} A dedicated check
confirms that an exploratory run covers the same 241 records and 237 components; that no
biological component straddles a training and test fold or two inner folds; that the split
table and its checksum are identical to those of the primary run; that the input
checksums, base seed and fold counts match; that tuning is nested; that models shared with
the primary run reproduce its predictions exactly; and that the primary run's own files
still match their recorded checksums.

\paragraph{Negative controls would not by themselves create inference.} With 1,000 draws,
the gene-set $p$-values that are currently empty would become available, and the observed
$\Delta\mathrm{MAE}$ could be placed descriptively against a permuted-outcome null. Tiers A
and B would stay closed, because that gate also requires a formally demonstrated prediction
gain, which needs the interval method that the deferred coverage simulation was designed to
validate. A permutation test would be a legitimate alternative route, but adopting it
requires a protocol amendment recorded \emph{before} the null distribution is inspected;
written afterwards, it would be a post-hoc test chosen in sight of the result and would have
to be reported as such.

\paragraph{The within-tissue correlation artefact is avoided.} None of these analyses
reports the pooled within-tissue rank correlation produced by the first version of that
table; the corrected statistic is undefined wherever a model predicts a constant within a
tissue (Section~\ref{sec:correction}).
